\section{Effective dimension and sample complexity of emergence}

\label{app:effective-dimension}

As discussed in the main text,  training dynamics is found in practice to often display long plateaus followed by abrupt transitions, with new directions in representation space \emph{emerging} sequentially \cite{wei2022emergent,raventos2023pretraining,arora2023theory,schaeffer2023are}. We return in this section to the emergence of learned features at each layer, and in particular, we now make the noise level \(\tau^k_\ell(n)\) in \eqref{eq:lofi-emergence-threshold} explicit.  Throughout the section our results apply conditioned on a fixed kernel $K_{\ell-1}(\bm x,\bm x')$ defined by the features $\bz_{\ell-1}$ of the previous layers.

Recall the residual candidate class
\[
  \mathcal{S}^\ell_{k}
  =
  \left\{
  \varphi\in\mathcal H_{\ell-1}:
  \norm{\varphi}_{\mathcal H_{\ell-1}}=1,\,
  \varphi \perp \varphi_1,\cdots, \varphi_{k-1}
  \right\}.
\]
Let \(T_{\ell-1}:L^2(P_x)\to L^2(P_x)\) be the integral operator of the current kernel,
\[
    (T_{\ell-1}f)(\bm x)
    =
    \int K_{\ell-1}(\bm x,\bm x')f(\bm x')\,dP_x(\bm x'),
\]
and let \(\Pi^\perp_k\) denote the projection orthogonal to the selected features \(\varphi_1,\dots,\varphi_{k-1}\). The residual covariance operator is
\[
    T_{\ell,k}
    :=
    \Pi^\perp_k T_{\ell-1}\Pi^\perp_k .
\]
For a resolution parameter \(r>0\), define the {\it residual effective dimension}:
\begin{definition}[Residual effective dimension]
    \label{def:lofi-effective-dimension}
Let  \(\{\lambda^{\ell,k}_j\}_{j\geq 1}\) denote the eigenvalues of \(T_{\ell,k}\). The \emph{effective residual dimension} at scale \(r\) is defined as
    \begin{equation}
    D^\ell_k(r)
    \coloneqq
    \operatorname{Tr}\!\left[
        T_{\ell,k}(T_{\ell,k}+r I)^{-1}
    \right]
    =
    \sum_{j\geq 1}
    \frac{\lambda^{\ell,k}_j}{\lambda^{\ell,k}_j+r},
    \label{eq:lofi-effective-dimension}
\end{equation}
\end{definition}
The above notion of effective dimension is standard in kernel regression
\cite{caponnetto2007optimal,rudi2017generalization}, and similar quantities
appear in local Rademacher analyses of kernel classes
\cite{mendelson2003performance,bartlett2005local}. \(D^\ell_k(r)\) can be
interpreted as counting the number of residual directions whose variance is
above the resolution \(r\).

\begin{assumption}[Residual eigenfunction hypercontractivity]
\label{ass:residual-eigenfunction-hypercontractivity}
Let \(\mathcal E_{\ell,k}\) denote the
\(L^2(P_x)\)-closed span of the residual eigenfunctions
\(\{\phi_j^{\ell,k}\}_{j\ge1}\). There exists \(H_\ell<\infty\) such that every
\(g\in\mathcal E_{\ell,k}\) satisfies
\[
    \norm{g}_{L^4(P_x)}
    \le
    H_\ell \norm{g}_{L^2(P_x)}.
\]
Equivalently, the same inequality holds for every \(L^2\)-convergent expansion
\(g=\sum_j c_j\phi_j^{\ell,k}\) with \(\sum_jc_j^2<\infty\).

\end{assumption}

Such an assumption holds, in particular, for the polynomial eigenbases that
appear in random-feature kernels such as Hermite polynomials and spherical harmonics (\cite{mei2022generalizationerror}).


\begin{theorem}[Emergence sample complexity]
\label{prop:emergence}
Suppose that Assumption~\ref{ass:residual-eigenfunction-hypercontractivity}
holds and the kernel \(K_{\ell-1}\) and labels \(y\) are uniformly bounded. Let
\[
    r_\ell^{k,\star}
    \coloneqq
    \argmax_{r:\, r \leq \lambda_1^{\ell,k}}
    r\sqrt{D^\ell_k(r)} .
\]
Here \(\lambda_1^{\ell,k}\) is the top residual eigenvalue, hence the largest
possible variance scale of a unit-RKHS candidate in \(\mathcal S^\ell_k\).
Then, for any \(\delta>0\), there is a constant
\(\widetilde C_{\ell,H}>0\) such that with probability at least \(1-\delta\),
\begin{equation}
    \begin{aligned}
     \tau^k_\ell(n)
     &\coloneqq
     \sup_{\varphi \in  \mathcal{S}^\ell_{k}}
    \left|
    \widehat c_{\ell,n}(\varphi)-c_\ell(\varphi)
    \right|
    \le
    \widetilde C_{\ell,H}\,
    r_\ell^{k,\star}
    \sqrt{\frac{D^\ell_k(r_\ell^{k,\star})}{n}} .
    \end{aligned}
\end{equation} This gives, up to polylogarithmic factors, the sample
scale
\[
    n_k
    \gtrsim
    \frac{
        (r_\ell^{(k)})^2 D^\ell_k(r_\ell^{(k)})
    }{
        (\rho_\ell^{(k)})^2
    },
\]
for recovery of $\varphi_k$ conditional on the recovery of previous features $\varphi_1,\cdots,\varphi_{k-1}$.
\end{theorem}

\begin{remark}\label{rem:order}
The order of maximizers of the noise scale $r\sqrt{D^\ell_k(r)}$ may not perfectly match the order of maximizers of the population correlation $\abs{\mathbb{E}[y \varphi(\bx)^2]}$. Hence, the next emergent feature after recovering $\varphi_1,\cdots,\varphi_{k-1}$ may not be $\varphi_{k}$. For simplicity, we neglect the correction in sample complexity due to such order switches.
\end{remark}
\begin{proof}

 For \(r>0\), let
\[
    \mathcal S^\ell_k(r)
    :=
    \{\varphi\in\mathcal S^\ell_k:\E[\varphi(\bm X)^2]\le r\},
\]

We proceed by decomposing the fluctuations into contributions from $ S^\ell_k(r)$ for different variance scales $r$. Define
\begin{equation}
    \tau^k_\ell(r,n)
    \coloneqq
    \sup_{\varphi \in \mathcal S^\ell_k(r)}
    \abs{\hat{\mathbb{E}}[y\varphi^2]-\mathbb{E}[y\varphi^2]} .
\end{equation}
We will show that:
\begin{equation}
    \tau^k_\ell(r,n)
    \lesssim
    \widetilde C_{\ell,H}\,
    r\sqrt{\frac{D^\ell_k(r)}{n}}
    \label{eq:lofi-noise-floor}
\end{equation}

The above bound follows through an extension of Theorem~2.1 in
\cite{mendelson2003performance} which bounds the rademacher complexity of $S^\ell_k(r)$. Concretely, Theorem 2.1 in \cite{mendelson2003performance} provides that for a unit ball of an RKHS of a bounded kernel with eigenvalues \(\{\lambda_j\}_{j\ge 1}\):
\[
    \mathbb E R_n
    \left\{
        f\in\mathcal H:\norm{f}_{\mathcal H}\le 1,\,
        \mathbb E f(\bm X)^2\le r
    \right\}
    \lesssim
    \left(
        \frac{1}{n}\sum_{j\ge 1}\min\{r,\lambda_j\}
    \right)^{1/2}.
\]
Since for every \(r>0\) and every
eigenvalue \(\lambda_j\),
\[
    \frac12 \min\{r,\lambda_j\}
    \le
    \frac{r\lambda_j}{r+\lambda_j}
    \le
    \min\{r,\lambda_j\},
\]
the same bound holds for $D^\ell_k(r)$ defined in Definition \ref{def:lofi-effective-dimension}.

We now adapt the proof of Theorem~2.1 of
\cite{mendelson2003performance} to the quadratic process. By the boundedness of the labels $y$, $\tau^k_\ell(r,n)$ can be bounded up to constants by the rademacher complexity of $\varphi^2$ over $S^\ell_k(r)$, defined as:

\begin{equation}
 \E_\epsilon[\sup_{\varphi\in\mathcal S^\ell_k(r)}
    \left|
        \sum_{i=1}^n\epsilon_i\varphi(\bm X_i)^2
        \right|],
\end{equation}
where $\epsilon \in \{\pm 1\}$ denote standard rademacher random variables.

Let
\(\{\phi_j^{\ell,k}\}_{j\ge1}\) denote the \(L^2(P_x)\)-orthonormal eigenfunctions
of \(T_{\ell,k}\) 
Any $\varphi$ with $\norm{\varphi}_{\mathcal{H}}=1$ can be expressed in the basis w.r.t \(\{\phi_j^{\ell,k}\}_{j\ge1}\)  as
\[
    \varphi_\beta(\bm x)
    =
    \sum_{j\ge1}
    \beta_j\sqrt{\lambda_j^{\ell,k}}\,
    \phi_j^{\ell,k}(\bm x),
    \qquad
    \sum_j\beta_j^2\le 1.
\]
The constraint $\E[\varphi(\bm X)^2]\le r$ then translates to
\(\sum_j\lambda_j^{\ell,k}\beta_j^2\le r\). Hence, the set of coefficients for $\varphi \in S^\ell_k(r)$ is given by:
\[
    I_r
    :=
    \left\{
        \beta:
        \sum_j\beta_j^2\le 1,\quad
        \sum_j\lambda_j^{\ell,k}\beta_j^2\le r
    \right\}.
\]
The proof of Theorem~2.1 in 
\cite{mendelson2003performance} relates the set $I_r$ to an ellipsoid $B_r$ through the following inclusions:
\[
    B_r\subset I_r\subset \sqrt2\,B_r,
    \qquad
    B_r
    :=
    \left\{\beta:\sum_j\mu_j(r)\beta_j^2\le1\right\},
    \qquad
    \mu_j(r)=\left(\min\{1,r/\lambda_j^{\ell,k}\}\right)^{-1}.
\]
Thus the supremum over \(I_r\) is bounded, up to an absolute constant, by the
supremum over \(B_r\). Setting \(\theta_j=\sqrt{\mu_j(r)}\,\beta_j\), 
\(B_r\) is re-parameterized as the Euclidean unit ball \(\norm{\theta}_2\le1\). Next, define
\[
    a_j(r) \coloneqq \frac{\lambda_j^{\ell,k}}{\mu_j(r)}
    =
    \min\{\lambda_j^{\ell,k},r\}
\]

The decomposition $\varphi_\beta(\bm x)
    =
    \sum_{j\ge1}
    \beta_j\sqrt{\lambda_j^{\ell,k}}\,
    \phi_j^{\ell,k}(\bm x)$ can then be re-expressed as:
    \begin{align*}
         &\sum_{j\ge1}
    \beta_j\sqrt{\lambda_j^{\ell,k}}\,
    \phi_j^{\ell,k}(\bm x)\\&= \sum_j
    \sqrt{\mu_j(r)}\,\beta_j \frac{1}{\sqrt{\mu_j(r)}}\phi_j^{\ell,k}(\bm x)\\&= \sum_j
            \theta_j\sqrt{a_j(r)}\phi_j^{\ell,k}(\bm x)
    \end{align*}

\begin{equation}\label{eq:radbound}
    \sup_{\beta\in I_r}
    \left|
        \sum_{i=1}^n
        \epsilon_i\varphi_\beta(\bm X_i)^2
    \right|
    \lesssim
    \sup_{\norm{\theta}_2\le1}
    \left|
        \sum_{i=1}^n
        \epsilon_i
        \left(
            \sum_j
            \theta_j\sqrt{a_j(r)}\phi_j^{\ell,k}(\bm X_i)
        \right)^2
    \right|.
\end{equation}
Define:
\[
    v_{\ell,k,r}(\bm x)
    :=
    \left(\sqrt{a_j(r)}\,\phi_j^{\ell,k}(\bm x)\right)_{j\ge1},
\]
then the RHS in Equation \ref{eq:radbound} can be expressed as:
\begin{equation}\label{eq:matconc}
    \left\|
        \sum_{i=1}^n
        \epsilon_i
        v_{\ell,k,r}(\bm X_i)\otimes
        v_{\ell,k,r}(\bm X_i)
    \right\|_{\mathrm{op}}.
\end{equation}

Let
\[
    \psi_{\ell,k}(r)^2
    :=
    \sum_j a_j(r)
    \asymp rD^\ell_k(r).
\]

We now bound Equation \ref{eq:matconc} through a matrix concentration argument which requires bounding the operator norm of the following fourth-moment matrix:
\[
    M_{\ell,k}(r)
    :=
        \E\!\left[
            \norm{v_{\ell,k,r}(\bm X)}_2^2
            v_{\ell,k,r}(\bm X)
            \otimes
            v_{\ell,k,r}(\bm X)
        \right].
\]
We bound $\norm{M_{\ell,k}(r)}$ through Assumption \ref{ass:residual-eigenfunction-hypercontractivity}. For every \(\norm{u}_2=1\),
write \(g_u(\bm x)=\langle u,v_{\ell,k,r}(\bm x)\rangle\). Then
\[
    \E[g_u(\bm X)^2]
    =
    \sum_j u_j^2a_j(r)
    \le r,
\]
since $a_j(r) \leq r$. Subsequently by Cauchy--Schwarz and hypercontractivity (Assumption \ref{ass:residual-eigenfunction-hypercontractivity}), we obtain:
\[
    \begin{aligned}
    \left\langle u,M_{\ell,k}(r)u\right\rangle
    &=
    \sum_j a_j(r)\,
    \E\!\left[
        \big(\phi_j^{\ell,k}(\bm X)\big)^2g_u(\bm X)^2
    \right]                                                   \\
    &\le
    \sum_j a_j(r)\,
    \norm{\phi_j^{\ell,k}}_{L^4(P_x)}^2
    \norm{g_u}_{L^4(P_x)}^2                                   \\
    &\le
    H_\ell^4
    \psi_{\ell,k}(r)^2
    \E[g_u(\bm X)^2]
    \le
    H_\ell^4 r\,\psi_{\ell,k}(r)^2 .
    \end{aligned}
\]
Therefore
\[
   \norm{M_{\ell,k}(r)}_2
    \lesssim
    H_\ell^4 r\,\psi_{\ell,k}(r)^2
    \asymp
    H_\ell^4 r^2D^\ell_k(r).
\]
Matrix Bernstein then gives, up to constants and polylogarithmic
factors:
\[
    \E_\epsilon
    \sup_{\varphi\in\mathcal S^\ell_k(r)}
    \left|
        \sum_{i=1}^n\epsilon_i\varphi(\bm X_i)^2
    \right|
    \lesssim
    H_\ell^2\sqrt{nr}\,\psi_{\ell,k}(r)
    \asymp
    H_\ell^2 n r\sqrt{\frac{D^\ell_k(r)}{n}}.
\]
By symmetrization and the boundedness of labels, this
yields \eqref{eq:lofi-noise-floor}. Taking the maximum over shells gives
the stated bound with \(r_\ell^{k,\star}\).
\end{proof}

\subsection{Estimating \texorpdfstring{$\tau^k_\ell(n)$}{tau-l-k(n)} in Feature space}
We discuss here how $\tau^k_\ell(n)$ can be estimated directly in the feature space (rather than in function space, as we did so far). 

Define the true label, weighted covariance:
\begin{equation}
    \bm C^{(\ell)}
    \coloneqq
    \mathbb{E}[
    y\,
    \bz_{\ell-1}(\bm x)
    {\bz}_{\ell-1}(\bm x)^\top] ,
\end{equation}
and the unweighted covariance:
\begin{equation}
    \bm \Sigma^{(\ell)}
    \coloneqq
    \mathbb{E}[
    \bz_{\ell-1}(\bm x)
    {\bz}_{\ell-1}(\bm x)^\top].
\end{equation}

To estimate the cutoff $n^k_\ell$ for the number of samples required for the recovery of $\phi_k$, we proceed as follows:
\begin{enumerate}[leftmargin=2em]
    \item Let $\bv_1,\dots \bv_{k-1}$ denote the top $k-1$ eigenvectors of $\bm C^{(\ell)}$. Compute the residual covariance:
    \begin{equation}
        \bm\Sigma^{(\ell)}_k \coloneqq (I-\bv_{k-1}\bv_{k-1}^\top) \bm\Sigma^{(\ell)}_{k-1}(I-\bv_{k-1}\bv_{k-1}^\top).
    \end{equation}

    \item Let $\{\lambda^{\ell,k}_j\}_{j=1}^\infty$ denote the eigenvalues of $\bm\Sigma^{(\ell)}_k$ . Then, for any $r \in \mathbb{R}$, define:
    \begin{equation}
       D^k_{\ell}(r) \coloneqq \sum_{j\geq 1} \frac{\lambda^{\ell,k}_j}{\lambda^{\ell,k}_j+r}.
    \end{equation}
    
    \item Set $r^k_\ell{=}\argmax_{r \leq \lambda_1^{\ell,k}} r \sqrt{D^k_{\ell}(r)}$.
    We predict the emergence of $\bv_{k}$ with samples $n^k_{\ell}$ satisfying:
    \begin{equation}\label{eq:real-data-emergence-estimate}
        n^k_{\ell} \sim  \left(\frac{r^k_\ell}{\rho^{(k)}_\ell}\right)^2 D^k_{\ell},
    \end{equation}
    where $\rho^{(k)}_\ell$ is the $k_{th}$ eigenvalue of $\bm C^{(\ell)}$.
\end{enumerate}



\subsection{Effective dimension in the Hermite toy model}
\label{app:effective-dimension-hermite}

The abstract quantity \(D^\ell_k(r)\) becomes concrete in the toy model Appendix~\ref{app:th_rf}. In this paragraph, \(k\) denotes the
Hermite degree used to define the first hidden representation \(h^{(1)}\), not
the index of the sequentially selected feature. Write the flattened degree-\(k\)
Hermite tensor as
\[
    H_k(\bm x)
    =
    \{H_\alpha(\bm x):|\alpha|=k\}
    \in \mathbb R^{D_k},
    \qquad
    D_k=\binom{d+k-1}{k}.
\]
For Gaussian inputs, these coordinates are orthonormal in \(L^2(P_x)\). Hence,
for any two coefficient vectors \(a,b\in\mathbb R^{D_k}\),
\[
    \mathbb E\!\left[
        \langle a,H_k(\bm X)\rangle
        \langle b,H_k(\bm X)\rangle
    \right]
    =
    \langle a,b\rangle_{\mathbb R^{D_k}} .
\]

Thus a degree-\(k\) candidate feature
\(\varphi_a(\bm x)=\langle a,H_k(\bm x)\rangle\) is literally a vector in a
\(D_k\)-dimensional Euclidean feature space as far as \(L^2(P_x)\) norms and
projections are concerned.

Equivalently, the degree-\(k\) Hermite kernel
\[
    K_{H,k}(\bm x,\bm x')
    =
    \langle H_k(\bm x),H_k(\bm x')\rangle
    =
    \sum_{|\alpha|=k} H_\alpha(\bm x)H_\alpha(\bm x')
\]
has integral operator
\[
    T_{H,k}f
    =
    \sum_{|\alpha|=k}
    \langle f,H_\alpha\rangle_{L^2(P_x)} H_\alpha .
\]
This is the orthogonal projector onto the degree-\(k\) Hermite chaos. Its only
nonzero eigenvalue is \(1\), with multiplicity \(D_k\). If the same block is
reached through a random-feature activation, as in the construction of
Subsection~\ref{gif:rf_maths}, the eigenvalue is multiplied by the squared
Hermite coefficient and by normalization constants, but the multiplicity is
still at most \(D_k\).

Therefore,
effective dimension relevant for the recovery of \(h^{(1)}\) is bounded by this
Hermite block rank. For a block eigenvalue \(a_k>0\),
\[
    D_{\mathrm{toy},1}^{(k)}(r)
    =
    \sum_{j=1}^{D_k}\frac{a_k}{a_k+r}
    =
    \frac{a_k}{a_k+r}D_k
    \le D_k
    =
    \binom{d+k-1}{k}
    =
    O(d^k),
\]
and for resolutions \(\lambda\lesssim a_k\) this is
\(\asymp D_k\). Hence, during first-level
feature recovery, the statistical fluctuation is governed by the ambient
degree-\(k\) Hermite space.

The planted first-level variables in Appendix~\ref{app:th_rf},
\[
    h_i^{(1)}(\bm x)
    =
    \langle A_i^{(1)},H_k(\bm x)\rangle,
    \qquad i=1,\ldots,d_1,
\]
form a \(d_1=d^\epsilon\) dimensional signal subspace inside this
\(D_k=O(d^k)\) dimensional Hermite block. The residual effective dimension
controls the size of the empirical noise over the full candidate block, while
the planted subspace controls the rank and strength of the population signal.
Concretely, by the definition of $y$, we have that:
\begin{equation}
     \mathbb{E}[y(h_i^{(1)}(\bm x))^2] \sim \frac{1}{\sqrt{d_1}}
\end{equation}


Substituting the above scaling into the sample complexity prediction prescribed by Theorem \ref{prop:emergence} gives the
\(d^{k+\epsilon}\) first-stage sample scale stated in
Appendix~\ref{app:th_rf}.



\subsection{Effective dimension in the multi-index models}


We now discuss how the sample complexity prediction in Proposition \ref{prop:emergence} recovers the rates in the setting of multi-index models with power-law dependence on the label. 

We follow here \cite{defilippis2026scaling,defilippis2026optimal} and consider hierarchical multi-index target \(y = f^*(\langle\bm{w}_1^*,\bm{x}\rangle,\dots,\langle\bm{w}_r^*,\bm{x}\rangle) + \xi\) with isotropic Gaussian inputs \(\bm{x}\sim\mathcal{N}(0,\bm{I}_d)\) and orthonormal teacher directions \(\{\bm{w}_i^*\}_{i=1}^r\). Because the input covariance is the identity \(D^\ell \asymp d\), at the slice-wise scale used in Proposition~\ref{prop:emergence}. The features at layer \(\ell\) along each candidate direction have unit variance under the Gaussian input, so the variance scale is order one, \(\tau = r_\ell^{(k)} = O(1)\). Finally, the power-law dependence of the label on the teacher coefficients translates into a power-law decay of the population correlations across the planted directions: the \(i\)-th spike satisfies \(\rho_\ell^{(i)} = O(i^{-\gamma})\), where \(\gamma>0\) is the exponent governing the label spectrum. Plugging \(D_k^\ell \asymp d\), \(\tau = O(1)\) and \(\rho^{(i)} = O(i^{-\gamma})\) into the recovery condition \(n \gg (r_\ell^{(k)})^2 D_k^\ell / (\rho_\ell^{(k)})^2\) of Proposition~\ref{prop:emergence} yields the sample-complexity threshold \(n_i \gtrsim d\, i^{2\gamma}\) for resolving the \(i\)-th spike, which matches the optimal scaling laws derived in~\cite{defilippis2026scaling,defilippis2026optimal}.
